\section{Complex Derivative}
\label{sec:derivata_complessa}

If we have a function $f\colon A\subset \CC \to \CC$, we can define the
(complex) derivative exactly as we did for the real derivative:
\[
  f'(x) = \lim_{h\to 0} \frac{f(x+h)-f(x)}{h}
\]
(if the limit exists finitely). Note that in this case $h\in \CC$, the
difference quotient is a complex division, and the limit is a complex limit.

For example, the function $f(z)=z$ is differentiable in the complex sense, and
its derivative is $f'(z) = 1$ because
\[
  \lim_{h\to 0} \frac{z+h-z}{h} = \lim_{h\to 0} \frac{h}{h}=1.
\]
The proofs concerning the derivative of the sum and the derivative of the
product are formally identical to those we have done for the real derivative. We
can thus state that if $f$ and $g$ are differentiable in the complex sense at a
point $z_0\in \CC$, it follows that
\begin{align*}
  (f+g)'(z_0) 
    &= f'(z_0) + g'(z_0), \\
  (f\cdot g)'(z_0) 
    &= f'(z_0)\cdot g(z_0) + f(z_0)\cdot g'(z_0).
\end{align*}
Thus, if $n\in \NN\setminus\ENCLOSE{0}$, we have
\[
  \enclose{z^n}' = n z^{n-1}.
\]
Obviously, the derivative of a constant is zero. The formula for the derivative
of the composite function is also valid and is proven by formally repeating the
same proof we have already seen:
\[
  \enclose{f(g(z))}'  = f'(g(z))\cdot g'(z).
\]
The formula for the derivative of the reciprocal is also obtained with the same
proof we have already seen in the real case:
\[
  \enclose{\frac 1 z}' = - \frac{1}{z^2}
\]
and consequently, the formula for the derivative of the quotient holds
\[
\enclose{\frac{f(z)}{g(z)}}' = \frac{f'(z)\cdot g(z) - f(z)\cdot g'(z)}{g^2(z)}.
\]

Thus, every polynomial function with coefficients in $\CC$ is differentiable in
the complex sense. The same holds for rational functions, i.e., ratios of
polynomials.

The exponential function is another very important example of a complex
differentiable function because it results in
\[
  \lim_{h\to 0} \frac{e^{z+h}-e^z}{h}
  = \lim_{h\to 0} e^z \cdot \frac{e^h-1}{h} = e^z
\]
by virtue of the well-known remarkable limit enjoyed by the exponential function
(Theorem~\ref{th:exp_complesso}). More generally, it is easy to verify that all
analytic functions are differentiable in the complex sense.

A surprising result of complex analysis tells us that the converse is also true:
every function differentiable in the complex sense at all points of its domain
is analytic (and in particular of class $C^\infty$).

In fact, differentiability in the complex sense is a very strong property. For
example, the function $f(z) = \bar z$ is not differentiable in the complex sense
because
\[
\frac{\overline{z+h}-\overline z}{h} = \frac{\bar h}{h}
\]
and for $h\to 0$, the limit does not exist since if $h$ is real, this ratio is
always $1$, but if $h$ is purely imaginary, this ratio is always $-1$.